\documentclass[10pt]{scrartcl}
\usepackage[mdthm,colorsec,noauthor,noasy]{evan}
\usepackage{handout}
\newcommand{\ip}[2]{\langle #1,#2\rangle}
\DeclareMathOperator{\tr}{tr}\DeclareMathOperator{\nul}{null}\DeclareMathOperator{\range}{range}\DeclareMathOperator{\spn}{span}
\newcommand{\LL}{\mathcal L}\newcommand{\PP}{\mathcal P}\newcommand{\MM}{\mathcal M}
\title{Math 110 --- Week 8}
\subtitle{Solutions}
\author{}\date{}\ihead{}
\begin{document}
\maketitle
\begin{prob}
\textbf{HW8 P1. }\par 
\begin{enumerate}[(a)]
  \ii \idea{Differentiation lowers degree, so enough derivatives kill any polynomial; the index is the first power that kills $x^n$, the slowest vector to die.}$D$ lowers degree by $1$ and kills constants, so $D^{n+1}p=0$ for every $p\in\PP_n(\CC)$, while $D^nx^n=n!\ne0$. So $D$ is nilpotent of index $n+1$.
  \ii $E(0,D)=\nul D=$ constants, so the geometric multiplicity is $1$. $G(0,D)=\PP_n(\CC)$ since $D$ is nilpotent, so the algebraic multiplicity is $n+1$.
  \ii \idea{One Jordan block means one chain. Start from the vector that takes longest to die ($x^n$) and differentiate repeatedly; dividing by factorials makes each step land on the next basis vector with coefficient $1$.}$f_k=x^k/k!$ for $k=0,\dots,n$: $Df_0=0$ and $Df_k=x^{k-1}/(k-1)!=f_{k-1}$. In the basis $f_0,f_1,\dots,f_n$, column $k$ of the matrix has a single $1$ in row $k-1$: the matrix is the single block $J_{n+1}(0)$.
\end{enumerate}
\end{prob}

\begin{prob}
\textbf{HW8 P2. }\par 
\begin{enumerate}[(a)]
  \ii \idea{Translate the data into blocks: geometric multiplicity $=$ number of blocks, algebraic $=$ total size, largest block $=$ exponent in $m$.}Number of blocks $=$ geometric multiplicity, total size $=$ algebraic multiplicity: $J_3(2)\oplus J_1(-1)\oplus J_1(-1)$,
  \[ \begin{psmallmatrix}2&1&0&0&0\\0&2&1&0&0\\0&0&2&0&0\\0&0&0&-1&0\\0&0&0&0&-1\end{psmallmatrix},\qquad m=(x-2)^3(x+1). \]
  \ii Two blocks for $2$ with total size $3$: $J_2(2)\oplus J_1(2)\oplus J_1(-1)\oplus J_1(-1)$, i.e.\ the matrix in (a) with the $(2,3)$ entry changed to $0$. The largest block for $2$ has size $2$, so $m=(x-2)^2(x+1)$.
\end{enumerate}
\end{prob}

\begin{prob}
\textbf{HW8 P3. }\par 
\begin{enumerate}[(a)]
  \ii $N=J-I$ has $Ne_1=0$, $Ne_2=e_1$, $Ne_3=e_2$. So $\nul N=\spn(e_1)$; $N^2e_3=e_1$ and $N^2e_1=N^2e_2=0$, so $\nul N^2=\spn(e_1,e_2)$; $N^3=0$, so $\nul N^3=\FF^3$.
  \ii $G(1,J)=\nul(J-I)^3=\FF^3$. Geometric multiplicity $\dim\nul(J-I)=1$; algebraic multiplicity $\dim G(1,J)=3$.
  \ii \idea{Split $J$ into a scalar part and a nilpotent part that commute, then use the binomial theorem; it stops at $N^2$ because $N^3=0$.}$J=I+N$ with $IN=NI$ and $N^3=0$, so $J^m=I+mN+\binom m2N^2=\begin{psmallmatrix}1&m&\binom m2\\0&1&m\\0&0&1\end{psmallmatrix}$ (with $\binom m2=\frac{m(m-1)}2$; for $m=1$ this is $J$).
\end{enumerate}
\end{prob}

\begin{prob}
\textbf{HW8 P4. }\par 
\begin{enumerate}[(a)]
  \ii \idea{Null spaces of a block diagonal matrix add block by block, and for one block $J_j(0)^k$ is a shift by $k$ places.}For one block $J_j(0)$, $\dim\nul J_j(0)^k=\min(k,j)$. Add over the blocks of sizes $3$ and $2$: $k=0$: $0$; $k=1$: $1+1=2$; $k=2$: $2+2=4$; $k=3$: $3+2=5$; $k=4$: $3+2=5$.
  \ii \idea{Differences of the $d_k$ count blocks of size $\ge k$; the total dimension pins down the largest size.}The number of blocks of size $\ge k$ is $\dim\nul N^k-\dim\nul N^{k-1}$: $2-0=2$ blocks of size $\ge1$, $4-2=2$ of size $\ge2$, $5-4=1$ of size $\ge3$. So there are two blocks, both of size at least $2$, exactly one of size at least $3$. The three numbers alone do not fix the larger size ($J_4(0)\oplus J_2(0)$ on $\CC^6$ also gives $2,4,5$), so we use that $N$ acts on the same $5$-dimensional space as in (a): then $\nul N^3=V$ and $N^3=0$, so no block has size above $3$. Hence the sizes are $3$ and $2$. (Equivalently: the sizes are $2$ and $s\ge3$, and $2+s=\dim V=5$.)
\end{enumerate}
\end{prob}

\begin{prob}
\textbf{HW8 P5. }\par 
\idea{For each eigenvalue, make the table of $d_k$ and take differences (Prop.~77).}Eigenvalue $1$: $d_1,d_2,d_3=1,2,2$, so $1$ block of size $\ge1$, $1$ of size $\ge2$, none of size $\ge3$: $J_2(1)$. Eigenvalue $2$: $d_1,\dots,d_4=2,3,4,4$, so $2$ blocks of size $\ge1$, $1$ of size $\ge2$, $1$ of size $\ge3$, none of size $\ge4$: $J_3(2)$ and $J_1(2)$. Check: $2+3+1=6$. So $T\sim J_2(1)\oplus J_3(2)\oplus J_1(2)$.
\end{prob}

\begin{prob}
\textbf{HW8 P6. }\par 
(a) $\chi=(x-4)^5(x+1)$. (b) $m=(x-4)^3(x+1)$ (the largest block for $4$ has size $3$). (c) Eigenvalue $4$: geometric $2$ (two blocks), algebraic $5$; eigenvalue $-1$: geometric $1$, algebraic $1$. (d) Not diagonalizable: $m$ has the repeated root $4$ (equivalently, $2\ne5$).
\end{prob}

\begin{prob}
\textbf{HW8 P7. }\par 
\idea{Everything about a block diagonal matrix splits into blocks, so count block by block; for one block the null space of $J_j(0)^k$ is easy because $J_j(0)^k$ is a shift.}Choose a basis in which $N$ is in Jordan form with blocks $J_{j_1}(0),\dots,J_{j_r}(0)$. $N^k$ is block diagonal with blocks $J_{j_i}(0)^k$, so $\dim\nul N^k=\sum_i\dim\nul J_{j_i}(0)^k$. On one block with basis $e_1,\dots,e_j$, $J_j(0)^ke_l=e_{l-k}$ for $l>k$ and $0$ for $l\le k$; so $\nul J_j(0)^k=\spn(e_1,\dots,e_{\min(k,j)})$ and $\dim\nul J_j(0)^k=\min(k,j)$. Then $\min(k,j)-\min(k-1,j)$ is $1$ if $j\ge k$ and $0$ if $j<k$. Summing over the blocks, $\dim\nul N^k-\dim\nul N^{k-1}$ counts the blocks with $j_i\ge k$.
\end{prob}

\begin{prob}
\textbf{G1. }\par 
\idea{Get the eigenvalues from a characteristic polynomial, then make the $d_k$ table, then build chains from the top.}
The upper-left block has trace $4$ and determinant $4$, so $\chi_A=(x-2)^2\cdot(x-2)=(x-2)^3$: only eigenvalue $2$. $A-2I=\begin{psmallmatrix}-1&1&0\\-1&1&0\\0&0&0\end{psmallmatrix}$ has rank $1$, so $E(2,A)=\{y=x\}=\spn\big((1,1,0),(0,0,1)\big)$ and $d_1=2$. $(A-2I)^2=0$, so $d_2=3$ and $G(2,A)=\CC^3$. Blocks: $2$ of size $\ge1$, $1$ of size $\ge2$: $J_2(2)\oplus J_1(2)$, $m_A=(x-2)^2$. Jordan basis: top $e_1\notin\nul(A-2I)$, bottom $(A-2I)e_1=(-1,-1,0)$; then the eigenvector $e_3$, independent of that bottom. In the basis $(-1,-1,0),e_1,e_3$ the matrix is $J_2(2)\oplus J_1(2)$ (check: $Ae_1=(1,-1,0)=2e_1+(-1,-1,0)$).
\end{prob}

\begin{prob}
\textbf{G2. }\par 
\idea{Same as for eigenspaces: $S$ commutes with every polynomial in $T$, so it preserves null spaces of polynomials in $T$.}
$G(\lambda,T)=\nul(T-\lambda I)^n$. If $(T-\lambda I)^nv=0$, then $(T-\lambda I)^nSv=S(T-\lambda I)^nv=0$.
\end{prob}

\begin{prob}
\textbf{G3. }\par 
\begin{enumerate}[(a)]
  \ii \idea{$(\lambda I+N)v=0$ says $v$ is an eigenvector of $N$ for $-\lambda$, and $0$ is the only eigenvalue of $N$.}
  If $\lambda\ne0$ and $(\lambda I+N)v=0$, then $Nv=-\lambda v$; since $-\lambda$ is not an eigenvalue of $N$, $v=0$. So $\lambda I+N$ is injective, hence invertible. If $\lambda=0$, $N$ is not injective ($V\ne0$ and $N^k=0$). (Explicitly, $(\lambda I+N)^{-1}=\lambda^{-1}\sum_{j=0}^{n-1}(-N/\lambda)^j$.)
  \ii \idea{With one eigenvalue there is one generalized eigenspace, and it is everything.}
  By Prop.~73, $V=G(\lambda,T)=\nul(T-\lambda I)^n$, so $(T-\lambda I)^n=0$.
\end{enumerate}
\end{prob}

\begin{prob}
\textbf{G4. }\par 
\idea{Block sizes add to the algebraic multiplicity $4$; the largest is the exponent $2$ in $m_T$.}
Sizes with largest $2$ and sum $4$: $\{2,2\}$ or $\{2,1,1\}$. So $J_2(2)\oplus J_2(2)$ or $J_2(2)\oplus J_1(2)\oplus J_1(2)$. They have the same $\chi_T$ and $m_T$; $\dim E(2,T)$ ($2$ versus $3$) tells them apart.
\end{prob}

\begin{prob}
\textbf{G5. }\par 
\idea{Over $\CC$, diagonalizable means $E(\lambda,T)=G(\lambda,T)$ for every $\lambda$.}
If $0\ne v\in G(\lambda,T)$, the hypothesis gives $Tv=\mu v$ for some $\mu$; then $0=(T-\lambda I)^nv=(\mu-\lambda)^nv$, so $\mu=\lambda$. Hence $G(\lambda,T)\subseteq E(\lambda,T)$, so $G(\lambda,T)=E(\lambda,T)$ for each eigenvalue. Then $V=\bigoplus G(\lambda_i,T)=\bigoplus E(\lambda_i,T)$ (Prop.~73), and a union of bases of the eigenspaces is an eigenbasis.
\end{prob}

\end{document}
